\documentclass[10pt]{article}
\usepackage[T1]{fontenc}
\usepackage[utf8]{inputenc}
\usepackage[french]{babel}
\usepackage{amsmath,amssymb,mathrsfs,mathpazo}
\usepackage{geometry}
\geometry{a4paper,margin=2cm}
\usepackage{graphicx,array,tabularx,tabularray,colortbl,multirow,multicol,booktabs,enumitem,pifont,ulem,fancybox,diagbox,cancel}
\usepackage{pstricks,pstricks-add,pst-plot,pst-tree,pst-node,pst-text,pst-func,pst-math,pst-bspline,pst-3dplot,pst-eucl,pst-solides3d}
\let\pscalculate\undefined
\usepackage{pst-bezier}
\usepackage{tikz,pgfplots,tkz-tab}
\usetikzlibrary{arrows,arrows.meta,patterns,calc,positioning,trees,decorations.markings,shapes.geometric,shapes.misc,shapes.arrows,intersections,decorations.pathmorphing,decorations.pathreplacing,backgrounds,shadows,babel}
\pgfplotsset{compat=1.18}
\usepgfplotslibrary{fillbetween}
\usepackage[np]{numprint}
\usepackage{listings,makecell,pas-tableur,twoopt}
\let\barre\undefined
\usepackage[pstricks]{tablvar}
\definecolor{darkgreen}{RGB}{0,100,0}
\providecommand{\red}{\color{red}}
\providecommand{\blue}{\color{blue}}
\providecommand{\green}{\color{green}}
\providecommand{\black}{\color{black}}
\providecommand{\white}{\color{white}}
\providecommand{\gray}{\color{gray}}
\providecommand{\cyan}{\color{cyan}}
\providecommand{\magenta}{\color{magenta}}
\providecommand{\yellow}{\color{yellow}}
\providecommand{\violet}{\color{violet}}
\providecommand{\orange}{\color{orange}}
\providecommand{\euro}{\texteuro}
\providecommand{\EUR}{\texteuro}
\providecommand{\ieme}{\textsuperscript{e}}
\providecommand{\degres}{\ensuremath{^\circ}}
\providecommand{\degree}{\ensuremath{^\circ}}
\providecommand{\vv}[1]{\overrightarrow{#1}}
\providecommand{\Coord}[1]{#1}
\providecommand{\decosix}{}
\providecommand{\starredbullet}{\ensuremath{\bullet}}
\pagestyle{empty}
\setlength{\parindent}{0pt}
\usepackage{adjustbox}\DeclareUnicodeCharacter{2236}{:}\DeclareUnicodeCharacter{2190}{\ensuremath{\leftarrow}}
\providecommand{\R}{\mathbb{R}}
\providecommand{\N}{\mathbb{N}}
\providecommand{\Z}{\mathbb{Z}}
\providecommand{\Q}{\mathbb{Q}}
\providecommand{\D}{\mathbb{D}}
\providecommand{\C}{\mathbb{C}}
\providecommand{\e}{{\mathrm{e}}}
\providecommand{\vect}[1]{\overrightarrow{#1}}
\providecommand{\vectt}[1]{\overrightarrow{\mathrm{#1}}}
\providecommand{\vvt}[1]{\overrightarrow{\mathrm{#1}}}
\providecommand{\np}[2][]{{#2~#1}}
\providecommand{\vv}[1]{{\overrightarrow{#1}}}
\renewcommand{\d}{{\mathrm{d}}}
\providecommand{\degres}{{^\circ}}
\providecommand{\degre}{{^\circ}}
\providecommand{\blue}{\color{blue}}
\providecommand{\red}{\color{red}}
\providecommand{\enskip}{\quad}
\providecommand{\llbracket}{\lbrack\!\lbrack}
\providecommand{\rrbracket}{\rbrack\!\rbrack}
\providecommand{\euro}{{\text{€}}}
\providecommand{\no}{{\text{n°}}}
\providecommand{\interval}[2]{{[#1;#2]}}
\providecommand{\intervFO}[2]{{[#1;#2[}}
\providecommand{\dotfill}{\ldots}

\providecommand{\ds}{\displaystyle}
\providecommand{\textding}[1]{#1}
\providecommand{\Oij}{\(\mathrm{O};\vec{\imath},\vec{\jmath}\)}
\providecommand{\Oijk}{\(\mathrm{O};\vec{\imath},\vec{\jmath},\vec{k}\)}
\providecommand{\Ouv}{\(\mathrm{O};\vec{u},\vec{v}\)}
\providecommand{\intervFF}[2]{\ensuremath{[#1;#2]}}
\renewcommand{\euro}{\ifmmode\text{\texteuro}\else\texteuro\fi}
\DeclareUnicodeCharacter{00B0}{\ifmmode^{\circ}\else\textdegree\fi}
\DeclareUnicodeCharacter{20AC}{\ifmmode\text{\texteuro}\else\texteuro\fi}
% Énoncés : Éducation nationale. Transcription et corrigés : APMEP, Denis Vergès et contributeurs. Découpage et classement : CoopMaths, Sébastien Lozano, Éric Elter et contributeurs. Adaptation de la présentation : Maths Explorer.
% CC BY-SA : https://coopmaths.fr/www/about/
% https://coopmaths.fr/alea/static/eam/2026/tex/eam_specifique_2026_06_polynesie_3_cor.tex
\begin{document}
\noindent\textit{APMEP / CoopMaths. CC BY-SA. Présentation adaptée par Maths Explorer.}\par\medskip
\begin{enumerate}
\item Par lecture graphique :
	\begin{enumerate}
		\item
		$f(5)\approx 180$. \quad La cinquième semaine le nombre de malades est proche de 180.

		\item
		Le pic est atteint pour $x=8$, donc lors de la huitième semaine de l'épidémie.

		Le nombre maximum de malades est entre 240 et 280, environ 260.
	\end{enumerate}

\item
	\begin{enumerate}
		\item
		 $f(1) = - 1^{3} + 12\times 1^{2} + 4 = -1 + 12 + 4 = 15 $.

		 Selon ce modèle, 15 personnes sont malades à la fin de la première semaine.

		\item Pour tout réel $x$ de l'intervalle $[0 ; 12]$,
		\quad$f'(x)= (-1) \times 3x^2 + 12\times 2x = -3x^2+24x$.

		et \quad$3x (-x+8) = -3 x^2 + 24x$.

		Ainsi $f'(x) = 3x (-x+8)$.

		\item
		$f(0) =-0^3 + 12\times 0^{2} + 4 = 4$

		$f(12) =-12^3 + 12\times 12^{2} + 4 = -12^3 + 12^3 + 4 = 4$

		Pour tout $x\in ]0 ; 12]$,\quad  $ 3x > 0 $

		$\begin{aligned}[t]
			-x+8 = 0 &\iff 8 = x \\
			&\iff x = 8
		\end{aligned}$
		\quad et \quad
		$\begin{aligned}[t]
			-x+8 > 0 &\iff 8 > x \\&\iff x < 8
		\end{aligned}$


		\begin{center}
			\begin{tikzpicture}
					\tkzTabInit[lgt=3.5,espcl=4]
					{	$x$/0.7,
						Signe de $3x$/1.4,
						Signe de $(-x+8)$/1.4,
						Signe de $f'(x)$/1.4,
						Variations de $f$/3.4
					}
					{$0$ ,8 , $12$}
				\tkzTabLine{z,+,t,+,}
				\tkzTabLine{t,+,z,-,}
				\tkzTabLine{z,+,z,-,}
				\tkzTabVar{-/4,+/$f(8)$,-/4}
			\end{tikzpicture}
		\end{center}


		\item Déterminer la valeur exacte du nombre de malades la semaine où a eu lieu le pic de l'épidémie.

		Le maximum est atteint pour $x=8$.

		$f(8) = -8^3 + 12 \times 8^2 + 4 = -512 + 768 + 4 = 256 + 4 = 260$

		Le maximum du nombre de malades est $f(8) = 260$

	\end{enumerate}
\end{enumerate}

\end{document}
